\section{Related Work}
\label{sec:related}

\subsection{Within-Year Seed Stability: B.7}

The companion paper B.7 \citep{b7paper} establishes seed stability for
minimum-edge-cut redistricting within a single census year.
Running 1000 different random initialisations of the METIS algorithm on
2020 Census data, B.7 finds that the natural ratio is consistent across seeds
for $>95\,\%$ of states and that partisan seat counts vary by at most one seat
in competitive states.
This establishes the MEC solution space as constrained: the algorithm converges
to a small number of topologically distinct solutions, with the minimum-edge-cut
solution dominating.

B.15 extends this stability analysis along a second dimension.
Where B.7 fixes the input data and varies the seed, B.15 fixes the algorithm
and varies the input data across three independent decennial surveys.
The two stability dimensions are complementary: a state with high seed stability
but low census stability has an algorithm that reliably finds its natural
partition given any one census snapshot, but whose natural partition changes as
demographics evolve.
A state with high census stability but low seed stability would have a durable
geographic structure that is difficult to reliably detect.
In practice, we expect seed stability and census stability to be correlated:
states with simple geographic structure (clear urban cores) are both easy to
detect (high seed stability) and geographically durable (high census stability).

\subsection{GeoSection and Isoperimetric Normalisation: B.8}

GeoSection \citep{b8paper} introduced the normalised edge-cut
$\mathrm{EC}/\sqrt{\min(i,k-i)}$ as the criterion for selecting the natural
first-level split ratio.
The key finding is that for states with compact, geographically isolated
metropolitan cores, GeoSection \emph{confirms} the asymmetric split that
unnormalised minimum-edge-cut would also select, while for states with
distributed population, it correctly selects a more balanced split.

For cross-census stability, GeoSection's normalisation provides a theoretical
stability guarantee: the denominator $\sqrt{\min(i,k-i)}$ is a function of
the \emph{number of districts}, not the population distribution.
If the state's seat count is unchanged between censuses (the common case), the
normalisation formula is identical; only the edge-cut numerator changes with
population.
We show in \S\ref{sec:theory} that the natural ratio is stable whenever the
population shift is small relative to the isoperimetric gap between the winning
and second-best ratios.

\subsection{Temporal Stability of Redistricting Maps: Literature}

Temporal stability of redistricting outcomes has received limited direct
treatment in the quantitative redistricting literature, but several threads
are relevant.

\paragraph{Population geography as a stable driver.}
\citet{rodden2019} documents the deep roots of the urban-rural political
divide in American geography.
The geographic sorting of Democratic voters into dense urban areas is not a
recent phenomenon: it has been building since at least the 1970s, accelerating
through the 1980s and 1990s, and is now a structural feature of American
political geography.
\citet{chen2013}'s analysis of ``unintentional gerrymandering'' shows that
compact, geography-only redistricting algorithms systematically produce
Republican-advantaged maps because of this geographic sorting---and that this
bias has been present across multiple redistricting cycles.
Both findings predict high cross-census stability for GeoSection: if the
underlying geographic structure is stable, the algorithm's output should be
stable.

\paragraph{MCMC ensemble stability.}
\citet{chikina2017} and \citet{deford2019} (GerryChain/ReCom) use Markov chain
Monte Carlo methods to sample from the space of valid redistricting plans and
assess whether an enacted plan is an outlier.
These methods have been applied primarily within a single census year; their
cross-census behaviour is not characterised.
StabilitySection provides a complementary analysis: instead of sampling the
space of plans within one year and measuring variance, we hold the algorithm
fixed and measure variance across years.

\paragraph{Partisan stability in the literature.}
\citet{mcghee2014} and \citet{stephanopoulos2015} introduced the efficiency
gap as a measure of partisan bias in redistricting.
Their analysis of historical congressional maps shows that efficiency gaps are
often stable across redistricting cycles within a state---a form of empirical
cross-cycle stability driven by population geography.
Our CSS provides a more fine-grained algorithmic measure of the same
underlying phenomenon.

\subsection{Census Data Changes Between Decennial Years}

Between each pair of consecutive decennial censuses, the Census Bureau makes
two types of changes relevant to redistricting:

\paragraph{Population counts.}
The total population and its geographic distribution change.
Between 2000 and 2020, the US population grew by $18.2\,\%$, from 281.4M to
331.4M.
Growth was highly uneven: the Sun Belt states (TX, AZ, GA, FL, NV) grew by
$30$--$60\,\%$; declining states (WV, MS, AL) contracted.
Population shifts within states are even larger: suburban expansion around major
metros, rural depopulation, and the concentration of population in a smaller
number of very large metro areas.

\paragraph{Tract boundary redesign.}
Between 2000 and 2010, the Census Bureau split, merged, and redrawn
approximately $15\,\%$ of census tracts.
A similar rate of redesign occurred between 2010 and 2020.
The adjacency graph $G$ used by GeoSection is therefore not identical across
years: some edges are added (split tracts become adjacent to new neighbors) and
some are removed (merged tracts eliminate intermediate adjacencies).
We control for this in our stability decomposition (\S\ref{sec:decomposition}).

\paragraph{Seat count changes.}
Congressional apportionment changes with population.
States that gain seats (TX: 32 in 2010 $\to$ 36 in 2020) cannot have
identical maps in the two years because the number of districts changes.
For these states, we measure bisection-tree structural similarity rather
than raw district assignment overlap.
States that maintain the same seat count across censuses (the majority of
states) allow direct comparison of assignments.

\subsection{Cross-Census Validity in Data Pipeline Research}

The companion paper C.2 (cross-census validation) examines whether the
computational infrastructure produces consistent results across census years.
It tests that the same algorithmic pipeline, applied to different census
datasets, produces well-formed outputs with valid population balance.
B.15's focus is complementary but distinct: we study the \emph{mathematical
invariance} of the algorithm's optimal answer, not the \emph{computational
reliability} of the pipeline.
C.2 is a prerequisite for B.15 (we need the three census years in the
pipeline), and B.15 provides the theoretical explanation for the empirical
findings C.3 reports on temporal district assignment drift.
